% Verbatim from /Users/alingalatan/Downloads/cubic_chaos_partial_results (2).zip:optimal_direction.tex (16 November 2026); GPT (a ChatGPT session of the user), relayed by the user.
% Intended filepath: optimal_direction.tex
\documentclass[11pt]{article}
\usepackage[margin=1in]{geometry}
\usepackage{amsmath,amssymb,amsthm}
\newtheorem{theorem}{Theorem}
\newtheorem{corollary}{Corollary}
\newcommand{\E}{\mathbb E}
\newcommand{\He}{\operatorname{He}}
\title{Optimal-direction reduction for real cubic Wiener chaos}
\author{}
\date{}
\begin{document}
\maketitle
\noindent\textbf{Status.} All statements below are rigorous. They give an exact
normal form, necessary spectral constraints, and a quantitative sufficient
region. They do not prove the unrestricted cubic Gaussian $L^1$ inequality,
nor exhibit a counterexample to it.

Write a real cubic chaos as
\[
 h(g)=\sum_{i,j,k=1}^n T_{ijk}:g_i g_j g_k:,
\]
where $T$ is a real symmetric $3$-tensor and colons denote Wick products.
Thus $\E h^2=6\|T\|_{\mathrm F}^2$.

\begin{theorem}[A maximizing direction kills the linear cross term]
Suppose $T\ne0$. Choose a unit vector $e$ such that
\[
 a=T(e,e,e)=\max_{\|u\|=1}|T(u,u,u)|>0.
\]
After rotating $e$ to the first coordinate, write $g=(X,Y)$. Then
\[
 h(X,Y)=a\He_3(X)+XQ(Y)+R(Y),
\]
where $Q(Y)=Y^{\mathsf T}BY-\operatorname{tr}B$ with
$B=3T(e,\cdot,\cdot)|_{e^\perp}$, and $R$ is pure cubic chaos on $e^\perp$.
In particular the term $(X^2-1)L(Y)$ vanishes. Moreover
\[
 -3aI\preceq B\preceq \frac{3a}{2}I.
\]
For every unit $v\perp e$, putting
$\mu=v^{\mathsf T}Bv$ and $r_v=T(v,v,v)$, one has the stronger coupling
constraint
\[
 \boxed{27a r_v^2\leq(3a+\mu)^2(3a-2\mu).}
\]
\end{theorem}

\begin{proof}
First variation of $T(u,u,u)$ on the sphere at $u=e$ gives
$T(e,e,v)=0$ for every $v\perp e$. The coefficient of
$(X^2-1)L(Y)$ in the Wick expansion is exactly
$L(Y)=3T(e,e,Y)$, so it vanishes.

Fix unit $v\perp e$ and put $b=T(e,v,v)$, $r=T(v,v,v)$.
On the circle in this plane,
\[
 T(e\cos\theta+v\sin\theta,e\cos\theta+v\sin\theta,e\cos\theta+v\sin\theta)
 =a\cos^3\theta+3b\cos\theta\sin^2\theta+r\sin^3\theta.
\]
The second derivative at zero is $-3a+6b\leq0$, hence $b\leq a/2$.
Evaluating at $u=(e\pm\sqrt3v)/2$ and averaging removes $r$ and gives
$|a+9b|\leq8a$, hence $b\geq-a$. Since $\mu=3b$, these are the stated
matrix inequalities.

For the sharper bound, divide by $a$ and write $\beta=b/a$, $\rho=r/a$.
The condition that the absolute value of the binary cubic be at most one
on its unit circle is exactly
\[
 |1+3\beta t^2|+|\rho|t^3\leq(1+t^2)^{3/2}\qquad(t>0).
\]
Thus
\[
 |\rho|\leq\inf_{t>0}
 \frac{(1+t^2)^{3/2}-|1+3\beta t^2|}{t^3}.
\]
For $-1\leq\beta\leq1/2$, the infimum is
$(1+\beta)\sqrt{1-2\beta}$. Here is a direct verification. When
$1+3\beta t^2\geq0$, the derivative of the displayed quotient has the
sign of $1+\beta t^2-\sqrt{1+t^2}$. When the sign is negative, its derivative
has the sign of $-1-\beta t^2-\sqrt{1+t^2}$. For $\beta\ne0$, the minimum
occurs at $t^2=(1-2\beta)/\beta^2$, on the appropriate branch, and substitution
gives the value stated above. For $\beta=0$ the infimum is $1$ at infinity;
the endpoint $\beta=1/2$ is obtained as $t$ tends to zero. These observations
also show sufficiency of this binary-plane constraint.

Consequently
\[
 r^2\leq\frac{(a+b)^2(a-2b)}a.
\]
Substitute $b=\mu/3$ and multiply by $27a$.
\end{proof}

\begin{corollary}[Explicit small-residual region]
In the normal form above set
\[
 q=\|Q\|_2,\qquad r=\|R\|_2,\qquad
 c=\sqrt{2/\pi},\qquad t=\sqrt{5a^2+q^2}.
\]
If
\[
 (1-c^2)r^2+2ctr\leq c^2a^2,
\]
or equivalently
\[
 r\leq\frac{c a^2}{\sqrt{t^2+(1-c^2)a^2}+t},
\]
then
\[
 \E|h|\leq c\|h\|_2.
\]
This result needs no prior bound on the cubic $R$.
If a bound $\E|R|\leq cr$ is independently available, the sufficient
condition improves to $2tr\leq a^2$.
\end{corollary}

\begin{proof}
The reflection-odd bound follows directly by changing the measure of $X$ to
$|X|d\gamma_1(X)/c$:
\[
 \E|a\He_3(X)+XQ(Y)|\leq c\sqrt{5a^2+q^2}=ct.
\]
Indeed $X^2$ has law $\chi^2_2$ under this changed measure, so
$\E(X^2-3)^2=5$, while $\E Q=0$. The triangle inequality and
Cauchy--Schwarz give $\E|h|\leq ct+r$.
Hermite orthogonality gives $\|h\|_2^2=6a^2+q^2+r^2=t^2+a^2+r^2$.
The claimed condition is exactly the result of squaring
$ct+r\leq c\sqrt{t^2+a^2+r^2}$. If $\E|R|\leq cr$ is available, instead
use $\E|h|\leq c(t+r)$ and square once more.
\end{proof}

\paragraph{Limit of this reduction.}
The direction maximizing the diagonal tensor contraction removes $L$ but
does not remove $R$, nor does it make $R$ independent of $Q$. The displayed
directional constraints control the leading homogeneous tensor of $R$;
they do not by themselves compare the Gaussian absolute moment of the sum
with its variance. In particular, this is not an induction proof of the
unrestricted inequality.
\end{document}
